\documentclass[11pt,a4paper]{article}
% ===================================================================
%  Galerie aller empfohlenen tikzpeople-Figuren und Optionen
%  Kompilieren: pdflatex galerie_tikzpeople.tex
%  (Die Figur „mexican“ fehlt absichtlich: sie bedient ein Stereotyp.)
% ===================================================================
\usepackage[T1]{fontenc}
\usepackage{lmodern}
\usepackage[ngerman]{babel}
\usepackage[a4paper,margin=2cm]{geometry}
\usepackage{tikzpeople}
\pagestyle{empty}
\setlength{\parindent}{0pt}

% \Figur[Optionen]{Name} – Figur mit Beschriftung
\newcommand\Figur[2][]{%
  \begin{minipage}[b]{2.2cm}\centering
    \begin{tikzpicture}\node[#2, minimum size=1.5cm, #1]{};\end{tikzpicture}\\
    \scriptsize\texttt{#2}\ifx\relax#1\relax\else\\\texttt{#1}\fi
  \end{minipage}\hspace{1mm}}

\begin{document}
{\Large\bfseries\sffamily Alle Figuren}\par\medskip
\foreach \f in {alice,bob,bride,builder,businessman,charlie,chef,conductor,
  cowboy,criminal,dave,devil,duck,graduate,groom,guard,jester,judge,maninblack,
  nun,nurse,physician,pilot,police,priest,sailor,santa,surgeon}{\Figur{\f}}

\bigskip
{\Large\bfseries\sffamily Optionen}\par\medskip
\Figur[mirrored]{alice}
\Figur[evil]{charlie}
\Figur[good]{dave}
\Figur[monitor]{bob}
\Figur[shield,sword]{guard}
\Figur[saturated]{alice}
\colorlet{haut}{brown!60}%
\Figur[shirt=red,skin=haut]{dave}
\end{document}
